\documentclass[11pt,a4paper]{article}
\usepackage[margin=25mm]{geometry}
\usepackage{amsmath,amssymb,amsthm,hyperref}
\setlength{\emergencystretch}{2em}
\newtheorem{lemma}{Lemma}
\newtheorem{theorem}[lemma]{Theorem}
\newcommand{\gen}[1]{\langle#1\rangle}
\newcommand{\diag}{\operatorname{diag}}
\title{Four pairs in the binary equality case}
\author{Human Compression project: a component of the L044 recovery}
\date{}
\begin{document}\maketitle

Let the natural alternating group $A_{2m}$ have a string C-marking
of rank $m$, with $m\ge8$ even. Let $H=G_i$ be a transitive
maximal omission preserving two blocks of size $m$, called rows,
and suppose every distinguished maximal omission $H_j$ is
intransitive. Choose a row-swapping mark $l$ of $H$ with two
noncommuting neighbours $r,s$ on opposite sides in the string.
Thus $r,s$ commute. Let $C$ contain the other $m-4$ marks.

We acquire the complete contracted-forest theorem. Its graph $J$
consists of all simple contracted edges of $r,s$ on the $l$-pairs.
It has rank three. The graph obtained by adjoining one selected
eligible edge for every $C$ label is connected, and every selected
$C$ edge is a bridge. A row-preserving $R$ mark is a double
transposition with two distinct contracted edges; a row-swapping
$R$ mark has a contracted triangle and equals $l$ elsewhere.
Contraction is only a graph operation; $r,s$ need not permute
the contracted pairs.

\begin{theorem}\label{main}
The nonisolated vertices of $J$ are exactly four pairs, and $J$
is connected on them. If both $r,s$ preserve the two rows, the
configuration is impossible.
\end{theorem}

\section*{The common graph reduction}
Write $Y$ for the set of nonisolated vertices of $J$, and put
$K=\gen{l,r,s}$. Each component of $J$ lifts to one $K$-orbit
of twice its size. Collapse these components in the graph with
the selected $C$ edges; the resulting graph is a tree.

Every $g\in C$ acts on pair indices and has the form
$\diag(\pi)l^\chi$, with one global row character
$\chi\in\{0,1\}$. Its additional projected edges lie in $Y$.
Indeed a connecting path for an additional edge can contain no
selected edge of another $C$ label, since this would bypass that
eligible edge. Thus the new edge either lies in one $J$ component,
or joins the same two components as its own selected edge. If
one of these components is an isolated outside vertex, the matching
condition on $\pi$ excludes a distinct additional edge: that
vertex is already used by the selected transposition.

Let $c_-,c_+$ denote the possible outer $C$ neighbours of $r,s$.
Every other $C$ mark centralizes $K$. Such a mark preserves each
nontrivial $J$ component. Components of sizes three and two give
different $K$-orbit sizes. In the case of three components of
size two, their marked actions differ: one has only $r$, one
has both $r,s$, and one has only $s$. They cannot be interchanged
by an equivariant permutation. Thus only $c_-,c_+$ can join
a nontrivial component to another component.

A mark commuting with $l,r$ preserves the contracted support of
$r$, because this is the union of the $\gen{l,r}$-orbits larger
than two. The same observation applies to $s$. An edge to an
outside pair must therefore start outside the support of the
$R$ mark with which its label commutes.

\begin{lemma}\label{four}
The graph $J$ is connected on four nonisolated vertices.
\end{lemma}
\begin{proof}
Rank three and the two individual local forms leave component
sizes $4$, $3+2$, or $2+2+2$ for the nontrivial components.
In the last case $|Y|=6$: each $R$ graph has two disjoint
edges, one shared between them. The three nontrivial components
are marked-inequivalent as above. Collapsing the outside selected
$C$ trees as well leaves at least four vertices to be joined
by just the two possible core-joining labels. This is impossible.

Suppose $|Y|=5$, with components $P,Q$ of sizes three and two.
Both $R$ marks preserve rows. If one swapped them, its triangle
would fill $P$ and it would equal $l$ on $Q$. A row-preserving
other mark would have just one actual transposition above $Q$,
which cannot commute with $l$ there. Two row-swapping triangles
have union rank two on the same support, rank at least four when
they meet in at most one vertex, or union on four vertices when
they meet in two. None gives the present case.

The outside vertices form one selected tree $O$: otherwise the
collapsed graph has at least four vertices and just two available
joining labels. Its two joins, labelled $c_-,c_+$, connect
$P,Q,O$ in a tree. If they are $P$--$O$ and $Q$--$O$, the
outside path between them is a noncommutation path. Adjacent
projected edges share an outside vertex; additional edges inside
$Y$ cannot cancel this noncommutation. Together with the path
$c_-,r,l,s,c_+$, this makes a cycle in the string, impossible.

Otherwise one join $g$ connects $P$ to $Q$. Suppose $g$ commutes
with $r$, so the other join $d$ commutes with $s$. The $r$ graph
cannot lie entirely in $P$: it would give a six-point
$\gen{l,r}$-orbit there and only two-point orbits over $Q$,
prohibiting $g$. Hence $r$ has one edge on $Q$ and one on two
vertices of $P$. Equivariance forces $g$ to carry the entire
four-point orbit over $Q$ onto that over those two vertices of
$P$. Its projected involution exchanges the corresponding two
pairs of vertices. It fixes the remaining vertex $v$ of $P$:
its matching has no available partner inside $P\cup Q$, and
its only selected edge already joins these components. Connectivity
of $P$ in $J$ puts $v$ in the support of $s$.

If $d$ joins $Q$ to $O$, its inside endpoint is moved by $g$.
Then $g,d$ do not commute, since $g$ fixes the outside endpoint
and $d$ has no other outside edge. If $d$ joins $P$ to $O$,
it cannot leave at $v$, because it commutes with $s$. Its endpoint
is again moved by $g$, giving the same contradiction. But the
two outer labels $c_-,c_+$ must commute. Thus $|Y|=5$ is
also excluded, proving the assertion.
\end{proof}

\begin{lemma}\label{tails}
After Lemma~\ref{four}, a $C$ mark whose selected edge lies
entirely outside $Y$ centralizes $K$.
\end{lemma}
\begin{proof}
Collapse $Y$ to one vertex. The selected edges form a tree;
each connected component on outside vertices therefore has
exactly one edge to $Y$, called its exit. Suppose an outside
selected label $c$ fails to commute with $K$. It must be the
outer neighbour of $r$ or of $s$. Follow its outside selected
tree towards its exit $e$. Consecutive edges along this path
have distinct labels and fail to commute. If both other endpoints
are outside, their projected transpositions share just one
outside vertex. If one is inside $Y$, evaluation at the outside
endpoint of the other edge still shows noncommutation. Additional
edges inside $Y$ cannot cancel either disagreement.

The exit label also fails to commute with $K$: a centralizing
permutation preserves its unique eight-point orbit and cannot
leave it. Thus the path just constructed joins two distinct
$C$ labels, each adjacent to an $R$ mark. They cannot both be
adjacent to the same $R$ mark, since that mark already has $l$
as its other neighbour. They attach to opposite ends of $r,l,s$,
and their outside path forms a cycle with that core path. The
noncommutation graph is a subgraph of a string, so this is
impossible.
\end{proof}

\section*{Both neighbours preserve rows}
Assume this from now on. Two distinct commuting double
transpositions have disjoint actual supports, the same four-point
support with different Klein matchings, or one common actual
transposition with two other mutually disjoint transpositions.
Indeed commutation preserves or exchanges their transposition
orbits, giving precisely these alternatives.

If their supports coincide with different matchings, all four
points lie in one row. A distribution of two points per row
admits only one row-preserving matching; odd row sizes cannot
support one. Both marks therefore use all four pairs of $Y$,
so neither outer $C$ neighbour can exit. This contradicts
transitivity.

Suppose their actual supports are disjoint. If $r$ uses four
pairs in one row, the remaining four points are in the other
row, so again both supports use all four pairs and no exit
exists. If $r$ uses four pairs with one transposition in each
row, the only matching on the remaining points is its row
mirror; the two contracted graphs coincide, giving rank two
instead of three. Finally, if $r$ uses three pairs, name it
\[
                 r=(a_0\ b_0)(a_1\ c_1).
\]
The available points force $s=(b_1\ d_1)(c_0\ d_0)$. Direct
multiplication gives
\[
                  (rl)^3=\diag((b\ c))l=(sl)^3.
\]
This element is outside $\gen l$, contradicting
$\gen{r,l}\cap\gen{l,s}=\gen l$.

In the common-edge case the three simple edges form a path.
A star would require three mutually disjoint actual edges at
one $l$-pair, which has only two points. A shared end edge or
a shared middle edge gives, up to row interchange, reversal
and naming,
\[
\begin{array}{ll}
\mathrm A:&r=(a_0\ b_0)(c_0\ d_0),\quad
            s=(b_1\ c_1)(c_0\ d_0),\\[2pt]
\mathrm B:&r=(a_0\ b_0)(a_1\ c_1),\quad
            s=(a_1\ c_1)(c_0\ d_0).
\end{array}                                                   \tag{1}
\]
The actual coloured $K$ graph is an eight-point path. Its edge
colours in the two cases are respectively
\[
 l,r,l,s,l,\{r,s\},l,\qquad
 l,r,l,\{r,s\},l,s,l.
\]
Neither path reversal preserves these colours. Hence the
centralizer of the $K$ action on these eight points is trivial.
By Lemma~\ref{tails}, every selected outside mark fixes the core
pointwise. Its global row character is zero and it has no extra
edge inside $Y$. It is precisely the plain doubled selected
transposition. The outside selected trees are paths: three
incident plain transpositions would give a noncommuting triangle.
An exit meets an endpoint of its outside path: two incident tail
marks would give it two noncommuting neighbours in addition to
its adjacent $R$ mark, impossible in a string.

\section*{The exits in the two forms}
\begin{lemma}\label{Aexit}
In form $\mathrm A$ there is exactly one exit, at $a$, and
it and every subsequent tail mark are plain doubled transpositions.
\end{lemma}
\begin{proof}
Only $c_-$ can exit, because $r$ uses all four pairs. Its exit
is at $a$, the only pair outside the contracted support of $s$.
The centralizer of the six-point alternating path for
$\gen{l,s}$ gives two possibilities, where $x$ is the first
outside index:
\[
            c=\diag((a\ x)),\qquad
            c=\diag((a\ x)(b\ d))l.
\]
There is a tail mark $t=\diag((x\ y))$, since the outside
tree has $m-4\ge4$ vertices. In the second possibility,
\[
             e=(ct)^3=\diag((b\ d))l=(sl)^3.
\]
Thus $e$ lies in both $\gen{l,r,s,c}$ and $\gen{r,s,c,t}$.
It does not belong to $Q=\gen{r,s,c}$: the set
$\{a_0,b_0,x_1,d_1\}$ is invariant under every generator of
$Q$, but $e$ sends $a_0$ to $a_1$ outside it. This contradicts
the required intersection being $Q$. Only the plain exit remains,
and Lemma~\ref{tails} supplies every plain tail mark.
\end{proof}

\begin{lemma}\label{Bexit}
In form $\mathrm B$ there is exactly one plain exit and a
plain doubled tail. Up to reversal, the exit is at $b$.
\end{lemma}
\begin{proof}
A left exit joins $b$ to an outside index $x$. Its possibilities,
again from the six-point alternating-path centralizer, are
\[
               c=\diag((b\ x)),\qquad
               c=\diag((b\ x)(a\ d))l.
\]
The right exit is symmetric. In the second possibility put
$u=cl=\diag((b\ x)(a\ d))$. Inside $\gen{c,r,l}$,
the element $(rr^u)^2$ is a nontrivial pure row-one three-cycle
on $\{a,c,d\}$. The projection of $\gen{r,r^l,u}$ to that
row is $S_5$ on $a,b,c,d,x$: it contains $(a\ c),(a\ b)$,
and their $u$-conjugates give $(d\ c),(d\ x)$. These
transpositions form a connected tree. Conjugating the pure
three-cycle therefore gives all three-cycles, hence pure $A_5$
on that row; conjugation by $l$ gives pure $A_5$ on the other.
Since $s$ differs from $r$ by the pure row-zero double
transposition $(a\ b)(c\ d)$, this puts $s$ in
$\gen{c,r,l}$, contrary to IP.

Both plain exits cannot coexist. Write them as
$c=\diag((b\ x))$ and $d'=\diag((d\ y))$. Here $x\ne y$,
since otherwise the two exits would not commute. Each of the
groups $\gen{c,r,l,s}$ and $\gen{r,l,s,d'}$ contains
$w=\diag((a\ b))$. For the first, $rr^c$ is a pure row-zero
three-cycle on $\{a,b,x\}$. The row-zero projection of
$\gen{r,s,r^l,s^l}$ is $S_4$ on the four core indices. Its
conjugates supply all cycles $(u\ v\ x)$, hence pure $A_5$;
conjugation by $l$ supplies the other row. Multiplying $r$ by
a suitable pure row-one three-cycle then gives $w$. For the
second group use $ss^{d'}$ on $\{c,d,y\}$ and the same argument.

However $w\notin K$. On each core row let $S_4$ act on its
three partitions into pairs,
\[
                  ab|cd,\qquad ac|bd,\qquad ad|bc.
\]
Let $p,q\in S_3$ be the distinct transpositions induced by
$(a\ b),(a\ c)$, and let $\phi$ be the automorphism exchanging
$p,q$. It satisfies $\phi^2=1$, since these transpositions
generate $S_3$. Both $r,s$ induce $(p,q)$ on the two rows,
whereas $l$ exchanges rows. The graph of $\phi$ in
$S_3\times S_3$ is invariant under this exchange because
$\phi^2=1$. Hence every row-preserving element of $K$ has
its two images in that graph. But $w$ has image $(p,p)$,
which is outside it. The required intersection of the two
end subgroups is $K$, a contradiction. Exactly one plain exit
remains, with its plain tail by Lemma~\ref{tails}.
\end{proof}

\section*{The omitted endpoint}
In both remaining forms the noncommutations of $H$ give one
connected path. Thus the missing parent label is an endpoint;
an interior omission would disconnect that path. We exclude
both orientations of the missing even involution $\sigma$.

In form $\mathrm A$, put $O=\{a\}$ together with the outside
indices. The $C$ marks generate diagonal $S_O$. If $\sigma$
is next to $s$, it centralizes $\gen{l,r,C}$. This group has
a large orbit of size $2(m-2)$ over $O\cup\{b\}$, and a
four-point orbit over $\{c,d\}$. Products of $C$-conjugates
of $r$ generate pure $A_{m-2}$ on the first row of the large
orbit. A centralizing permutation preserves that support and
fixes it pointwise, since the centralizer of a natural alternating
group of degree at least six is trivial. Transitivity then
forces it to fix the whole large orbit. On the remaining four
points $\gen{l,r}$ is the natural $D_8$, whose centralizer is
$\{1,\diag((c\ d))\}$. Both elements preserve the original
binary rows. Hence the missing generator would preserve those
blocks, contradicting the natural alternating parent.

If $\sigma$ is next to the far tail, it centralizes the subgroup
with the last tail mark removed. There is one orbit of size
$2(m-1)$ and the last isolated $l$-pair. Put $O'=O\setminus
\{\text{last index}\}$. Products of conjugates of $r$ under
the remaining $C$ marks give pure $A_{m-3}$ on $O'\cup\{b\}$
in row zero. Here $m-3\ge5$, so the same centralizer argument
fixes the whole large orbit. The only remaining nonidentity
action is the odd transposition on the last pair, impossible.

In form $\mathrm B$, orient the unique exit at $b$ and put
$O=\{b\}$ together with the outside indices. At the endpoint
next to $s$, the commuting subgroup $\gen{l,r,C}$ has a large
orbit of size $2(m-1)$ and the isolated pair over $d$. Products
of $C$-conjugates of $r$ give pure $A_{m-2}$ on $O\cup\{a\}$
in row zero. Its support and transitivity force the large orbit
to be fixed pointwise; only an odd flip of the isolated pair
remains. At the other endpoint, use pure $A_{m-3}$ on
$O'\cup\{a\}$ instead. Again the large orbit of size $2(m-1)$
is fixed and only the odd flip on the last pair remains.
Every orientation is excluded, proving Theorem~\ref{main}.

\paragraph{Attribution and scope.}
The historical target and fracture strategy are from
Cameron--Fernandes--Leemans--Mixer, Proposition~3.5. The root
investigator supplied the common support reduction and the
classification, intersection and endpoint arguments in this paper.
The continuation investigator prepared this manuscript and added
the explicit attachment interface in Lemma~\ref{tails}; this
addition requires its own distinct check. The separate component
with a row-swapping neighbour uses a new central pairing. No
finite diagnostic is a premise, and no whole-source recovery is
claimed here.
\end{document}
